\documentclass[11pt]{article}
\usepackage[margin=1in]{geometry}
\usepackage{amsmath,amssymb,amsthm}
\usepackage{algorithm}
\usepackage{algpseudocode}
\usepackage{booktabs}
\usepackage{hyperref}
\usepackage{tikz}
\usepackage{xcolor}
\usetikzlibrary{arrows.meta,positioning,fit,backgrounds,shapes.geometric,calc}

\newtheorem{theorem}{Theorem}
\newtheorem{proposition}{Proposition}
\newtheorem{lemma}{Lemma}
\newtheorem{definition}{Definition}

\newcommand{\lo}{\underline}
\newcommand{\hi}{\overline}
\newcommand{\indep}{\perp\!\!\!\perp}
\newcommand{\DSE}{\mathcal{D}_{\mathrm{SE}}}
\newcommand{\DLCN}{\mathcal{D}_{\mathrm{LCN}}}
\newcommand{\CVE}{\texttt{CredalVE}}

\title{Exactness of Credal Variable Elimination on Chain-Graph LCNs\\
\large What CVE actually bounds, when the per-target bucket schedule is exact
across chains, trees, polytrees, DAGs and chain graphs, and how the tightening
schemes (D1--D5) act on it}
\author{IBM Research}
\date{}

\begin{document}
\maketitle

\begin{abstract}
Credal Variable Elimination (CVE, \texttt{lcn/inference/marginal/cn/cve.py}) is the hub of the
credal-network family: the per-target bucket-elimination engine against which Credal Cluster Tree
Elimination (CCTE), ApproxLP, Interval BP and Credal IJGP are all measured. This note pins down
three things about CVE against the implementation. \textbf{(i)~What it bounds.} CVE optimizes over
the \emph{strong extension} $\DSE$ of the compiled credal network, not the LCN's distribution set
$\DLCN$. More precisely, CVE does not merely return \emph{a} valid outer bound over $\DSE$: it
returns \emph{exactly} the strong-extension interval $[\lo q_{\mathrm{SE}},\hi q_{\mathrm{SE}}]$,
because it recomputes an elimination order per query with the target eliminated \emph{last}, so the
extreme-vertex combination that attains a bound survives to the final bucket. Since
$\DLCN\subseteq\DSE$, the reported interval is always a valid outer bound on $\DLCN$, and it is the
\emph{tightest} such bound any strong-extension propagation can give.
\textbf{(ii)~Which tightening schemes apply.} The five schemes of the companion note
(\texttt{docs/tighter\_approximation.tex}) split for CVE exactly as for CCTE: D1/D2/D3 are
\emph{build-stage} knobs on \texttt{CredalNetworkVertices.from\_lcn} that CVE inherits for
\emph{free} (it merely consumes tighter vertices, and being exact over $\DSE$ it reports the
tighter strong extension verbatim); D4 (\texttt{coupling="cross-family"}) is \emph{implemented
inside} CVE as a prune-time function filter \emph{together with} an order augmentation that forces
the coupled atoms into a common bucket; D5 is a \emph{separate} engine (\texttt{CredalJT}), not a
CVE knob. \textbf{(iii)~When it is exact.} Because CVE is exact over $\DSE$, its exactness with
respect to $\DLCN$ reduces to a single geometric question: \emph{when does $\DSE=\DLCN$?} We show
this holds on a clean singly-connected, cross-family-\emph{sentence}-free LCN --- Markov chain,
tree, polytree, and chain graph with compound nodes (\texttt{alarm}) --- so CVE is exact there; and
we exhibit the precise boundary --- a cross-family \emph{sentence} (\texttt{d4\_biting}) --- where
$\DSE\supsetneq\DLCN$ and CVE's \texttt{coupling="off"} interval, though still a valid and exact
$\DSE$ bound, is looser than $\DLCN$. On that instance the D4 order augmentation lets CVE reach the
exact bound for the coupled \emph{conditional} query while remaining sound, and CredalJT (D5) is the
reliable exact engine in general. This last point contrasts sharply with CCTE, whose single shared
order makes even the \emph{unconditional} bound fall strictly \emph{inside} $[\lo q_{\mathrm{SE}},\hi
q_{\mathrm{SE}}]$ on the same instance.
\end{abstract}

\tableofcontents

\section{Background: the CVE algorithm}\label{sec:bg}

\subsection{LCNs, the credal network, and the strong extension}
An LCN over binary atoms $\mathbf V=\{V_1,\dots,V_n\}$ is a set of probability \emph{sentences},
each Type-1 $\ell\le P(\varphi)\le u$ or Type-2 $\ell\le P(\varphi\mid\psi)\le u$, over
propositional formulas. Together with the conditional independencies of the \emph{Local Markov
Condition} (LMC) of the primal graph, the sentences define a set $\DLCN\subseteq\Delta(2^{\mathbf
V})$ of joint distributions. The marginal task for a query atom $a$ is $\lo P(a)=\min_{p\in\DLCN}
P(a)$ and $\hi P(a)=\max_{p\in\DLCN} P(a)$; exact inference (\texttt{ExactInference},
\texttt{exact.py}) solves one nonconvex NLP over the full $2^n$ joint per bound.

The \texttt{cn/} pipeline instead compiles the chain-graph LCN into a \emph{credal network}: a
directed graph over the families $\mathrm{ch}_v\mid\mathrm{pa}_v$ of the simplified structure graph,
whose local credal sets $K_v=K(\mathrm{ch}_v\mid\mathrm{pa}_v)$ are interval-valued and computed per
family. \texttt{CredalNetworkVertices.from\_lcn} enumerates the extreme points of every $K_v$
(\texttt{cnv.extreme\_points}), and \emph{all} the credal-network engines --- CVE, CCTE,
Interval~BP, ApproxLP and Credal~IJGP --- consume exactly that object. What they bound is the
\emph{strong extension}
\begin{equation}\label{eq:se}
  \DSE=\Bigl\{\textstyle\prod_v P_v \;:\; P_v\in K_v \text{ chosen independently per parent-config}\Bigr\},
\end{equation}
the set of joints obtained by letting each family, and each parent configuration of each family,
pick one extreme point of its local set freely and independently. Because the chain-graph
factorization $\prod_v P(\mathrm{ch}_v\mid\mathrm{pa}_v)$ reproduces the LMC exactly (this is what
\texttt{verify\_lmc\_factorization.py} certifies), every element of $\DSE$ satisfies every LMC
assertion. Hence
\begin{equation}\label{eq:incl}
  \DLCN\;\subseteq\;\DSE,
  \qquad\Longrightarrow\qquad
  [\lo q_{\mathrm{LCN}},\hi q_{\mathrm{LCN}}]\;\subseteq\;[\lo q_{\mathrm{SE}},\hi q_{\mathrm{SE}}],
\end{equation}
so any engine that brackets $\DSE$ produces a valid \emph{outer} bound on $\DLCN$. The inclusion in
\eqref{eq:incl} is strict exactly when a \emph{cross-family sentence} (a sentence whose atom scope
is contained in no single family scope) couples choices that the free product of \eqref{eq:se}
decouples. Sections~\ref{sec:exact} and~\ref{sec:boundary} turn on this distinction.

\subsection{The per-target bucket schedule}\label{sec:schedule}
\CVE{}'s public \texttt{run(evidence, elim\_heuristic, epsilon, coupling, verbosity)}
(\texttt{cve.py:84--165}) computes the bounds of \emph{every} non-evidence singleton atom by
\emph{looping} a single-query bucket elimination over the credal-network nodes. For each target node
it calls \texttt{\_run\_single\_query} (\texttt{cve.py:290--482}), which:
\begin{enumerate}
  \item \textbf{Builds initial potentials from the extreme points} (\texttt{cve.py:320--381}). Each
    family becomes a set of CPT-shaped functions: one function per Cartesian combination of a chosen
    extreme point in every parent configuration. A \texttt{Potential} is a finite \emph{set} of
    non-negative arrays over a shared scope; this set is precisely an explicit vertex enumeration of
    the family's contribution to $\DSE$.
  \item \textbf{Chooses an elimination order in which the target is last}
    (\texttt{cve.py:391--410}). With \texttt{elim\_heuristic="topological"} it takes the DAG
    topological order and removes the query; with \texttt{"min-fill"} it runs the greedy
    interaction-graph heuristic \texttt{min\_fill\_order(scopes, exclude=\{query\})}. Either way the
    query node's own bucket \emph{survives} to the end.
  \item \textbf{Runs one bucket-elimination pass} (\texttt{cve.py:421--455}): for each variable in
    the order, collect the potentials mentioning it, \emph{combine} them (elementwise product,
    scope union), \emph{marginalize} the variable out, and \emph{prune} dominated functions
    ($g$ dominates $f$ iff $g\ge f$ everywhere and $g>f$ somewhere; \texttt{potentials.py:124--148}).
  \item \textbf{Reads off the target bounds} (\texttt{cve.py:469--481}): every surviving function
    $f$ over the query scope is normalized to $f/\sum f$, and the per-state min/max across all
    surviving functions gives $[\lo P(\text{query}),\hi P(\text{query})]$.
\end{enumerate}
The two decisive features are that (a) each potential is an explicit set of strong-extension
vertices, and (b) the query is eliminated \emph{last}. Feature~(a) is what makes CVE a
\emph{strong-extension} method; feature~(b) is what makes it \emph{exact} over $\DSE$, and is
exactly what CCTE's single shared two-pass order gives up (\texttt{docs/ccte\_exactness.tex}). An
optional $\epsilon>0$ activates \texttt{\_run\_single\_query\_approx} (\texttt{cve.py:484}), an
FPTAS that permits $\epsilon$-domination in pruning (Mau\'a et al.\ 2012); with $\epsilon=0$ the
pruning is exact.

\section{What CVE bounds: exact over the strong extension}\label{sec:bounds}

\begin{proposition}[$\DSE$ target]\label{prop:se}
With \texttt{coupling="off"} and exact pruning ($\epsilon=0$), \CVE{} computes, for every query atom
$a$, exactly the strong-extension interval
$[\lo q_{\mathrm{SE}}(a),\hi q_{\mathrm{SE}}(a)]$
$=[\min_{P\in\DSE}P(a),\ \max_{P\in\DSE}P(a)]$. In particular it is always a valid \emph{outer}
bound on $\DLCN$.
\end{proposition}

\begin{proof}[Proof sketch]
\emph{Every function CVE forms is a strong-extension point.} The initial potentials
(\texttt{cve.py:320--381}) are, by construction, the extreme points of the families, i.e.\ the
generators of $\DSE$ in \eqref{eq:se}. The two potential operations preserve this: \texttt{combine}
takes an elementwise product of one function from each argument (a product of local choices, still a
point of the joint over the union scope), and \texttt{marginalize} sums a variable out (the marginal
of such a product). Pruning removes only \emph{dominated} functions, which can never attain a
componentwise extremum, so it changes no bound. Hence at the final step every surviving function
over the query scope is the query-marginal of some $P\in\DSE$, and the min/max read-off
(\texttt{cve.py:469--481}) ranges over exactly $\{P(a):P\in\DSE\}$ at its extreme vertices.

\emph{The query-last order loses no extremum.} Because $\DSE$ is a product of finitely many
polytopes, $\max_{P\in\DSE}P(a)$ is attained at a vertex of that product, i.e.\ at a specific
combination of family extreme points. Eliminating the query \emph{last} means that combination is
never marginalized away before the read-off, and it is never pruned, because at the query scope it
is (by definition of the extremum) not dominated. Thus the read-off sees it, and CVE returns the
exact $\DSE$ bound. \eqref{eq:incl} then gives the outer-bound guarantee on $\DLCN$.
\end{proof}

\noindent This is the sense in which CVE is the ``gold'' member of the family
(\texttt{docs/tighter\_approximation.tex}, comparison table): it is exact over $\DSE$, whereas
ApproxLP is \emph{inner}, Interval~BP and IJGP carry loopy outer slack, and CCTE is exact over $\DSE$
only when its single shared order happens to coincide with a query-last order --- which fails under
cross-family coupling (Section~\ref{sec:boundary}).

\section{When CVE is exact with respect to $\DLCN$}\label{sec:exact}

By Proposition~\ref{prop:se}, CVE returns $[\lo q_{\mathrm{SE}},\hi q_{\mathrm{SE}}]$ \emph{exactly}.
Its exactness with respect to the LCN --- $[\lo q_{\mathrm{SE}},\hi q_{\mathrm{SE}}]=[\lo
q_{\mathrm{LCN}},\hi q_{\mathrm{LCN}}]$ --- therefore reduces entirely to the geometric question of
\textbf{when $\DSE=\DLCN$}. The following proposition characterizes the clean case; the topology
enters only through singly-connectedness of the primal graph.

\begin{proposition}[Exactness on clean singly-connected LCNs]\label{prop:exact}
Let $\mathcal L$ be a chain-graph LCN whose primal graph is singly connected (Markov chain, rooted
tree, polytree, or a chain graph whose compound nodes leave the family DAG singly connected, e.g.\
\texttt{alarm}) and which has \emph{no cross-family sentence} (every sentence's atom scope fits
inside one family; only the LMC assertions may be cross-family). Then with \texttt{coupling="off"}
and $\epsilon=0$, \CVE{} returns the exact LCN marginal bounds for every atom.
\end{proposition}

\begin{proof}[Proof sketch]
There are two gaps between the CVE output and the LCN truth, and both close.

\emph{(1) $\DSE=\DLCN$ (no gap B).} A chain-graph product $\prod_v P(\mathrm{ch}_v\mid\mathrm{pa}_v)$
satisfies every LMC assertion by construction, so a free product in \eqref{eq:se} can leave $\DLCN$
\emph{only} by violating a \emph{sentence}. When every sentence is intra-family, each family's local
credal set $K_v$ already encodes every sentence constraining its scope, and the free product of
compliant families is again compliant. Hence $\DSE=\DLCN$; equivalently, there is no cross-family
coupling for the free product of \eqref{eq:se} to break. (This is precisely the content certified by
\texttt{verify\_lmc\_factorization.py} on chain/tree/polytree instances.)

\emph{(2) CVE is exact over $\DSE$ (no schedule loss).} This is Proposition~\ref{prop:se}: the
query-last per-target order guarantees the extreme vertex combination survives to the read-off.
Singly-connectedness is not needed for this half --- CVE's per-target order is exact over $\DSE$ on
\emph{any} topology --- but it is what makes half~(1) hold, so both halves together give exactness
w.r.t.\ $\DLCN$.

Combining, $[\lo q_{\mathrm{SE}},\hi q_{\mathrm{SE}}]=[\lo q_{\mathrm{LCN}},\hi q_{\mathrm{LCN}}]$
and CVE reports it exactly.
\end{proof}

\subsection{Topology by topology}\label{sec:topo}

\paragraph{Markov chain.} The primal graph is a path $x_0\!-\!x_1\!-\!\cdots\!-\!x_n$; every family is
$P(x_{i+1}\mid x_i)$ with a single parent. Singly connected, no cross-family sentence, so
$\DSE=\DLCN$ and CVE is exact. On \texttt{examples/chain.lcn} ($n=8$) CVE returns
$P(x_0)=[0.300,0.500]$, $P(x_2)=[0.328,0.640]$, $P(x_7)=[0.334,0.667]$, matching
\texttt{ExactInference} and CredalJT (cf.\ \texttt{docs/tighter\_approximation.tex},
\S\,chain, and \texttt{docs/ccte\_exactness.tex}).

\paragraph{Tree.} A rooted branching structure in which every node has at most one parent. Same
argument: single-parent families, singly connected, no cross-family sentence $\Rightarrow$
$\DSE=\DLCN$ $\Rightarrow$ exact. Confirmed on \texttt{examples/tree.lcn}.

\paragraph{Polytree.} Singly connected but with colliders (a node with several parents). The
families are now $P(\mathrm{ch}\mid\mathrm{pa}_1,\dots,\mathrm{pa}_k)$, and moralizing the colliders
induces co-parent independencies in the LMC. As long as every \emph{sentence} is intra-family,
$\DSE=\DLCN$ still holds (the free product respects the moralized LMC), so CVE is exact. Confirmed on
\texttt{examples/polytree.lcn}. (This is the same boundary ARIEL sits on ---
\texttt{docs/ariel\_exactness.tex} --- exact on singly-connected, and the LMC is preserved
structurally along the unique family path.)

\paragraph{Chain graph with compound nodes.} \texttt{examples/alarm.lcn} has five atoms and a
compound node $C\text{-}D$; the families are $P(A\mid B,E)$, $P(B)$, $P(E)$, $P(C\text{-}D\mid A)$.
The two large-scope LMC assertions $(C\indep B,E\mid D,A)$ and $(D\indep B,E\mid C,A)$ are
cross-family, but they are \emph{LMC assertions, not sentences}: the factorization reproduces them
for free. There is no cross-family \emph{sentence}, so $\DSE=\DLCN$ and CVE is exact, returning e.g.\
$P(A)=[0.087,0.100]$, $P(B)=[0.100,0.123]$, $P(E)=[0.050,0.074]$ (matching the brute-force oracle,
and more reliable than raw \texttt{ExactInference}, which can locally over-tighten ---
\texttt{docs/tighter\_approximation.tex}, alarm).

\paragraph{General / loopy DAG, or any cross-family \emph{sentence}.} When the primal graph has more
than one path between two atoms (loopy), or a sentence couples atoms across families, $\DSE\supsetneq
\DLCN$ can be strict: the free product of \eqref{eq:se} realizes joints that violate the coupling.
CVE still returns the \emph{exact} $\DSE$ interval (Proposition~\ref{prop:se}), hence a valid outer
bound, but it is no longer tight. This is the boundary of Section~\ref{sec:boundary}, and the regime
where the coupling schemes D4/D5 earn their keep.

\begin{center}\small
\begin{tabular}{lccc}
\toprule
Topology & singly conn.? & cross-family \emph{sentence}? & CVE (\texttt{coupling="off"}) \\
\midrule
Markov chain (\texttt{chain})     & yes & no  & \textbf{exact} \\
Tree (\texttt{tree})              & yes & no  & \textbf{exact} \\
Polytree (\texttt{polytree})      & yes & no  & \textbf{exact} \\
Chain graph, compound (\texttt{alarm}) & yes & no  & \textbf{exact} \\
Cross-family sentence (\texttt{d4\_biting}) & yes & \textbf{yes} & exact over $\DSE$; outer on $\DLCN$ \\
Loopy DAG                          & no  & --- & exact over $\DSE$; outer on $\DLCN$ \\
\bottomrule
\end{tabular}
\end{center}

\section{The boundary: a cross-family sentence (\texttt{d4\_biting})}\label{sec:boundary}

The minimal witness that separates $\DSE$ from $\DLCN$ is \texttt{examples/d4\_biting.lcn}: three
atoms, $A$ the root with independent children $B,C$, so the families are $P(A)$, $P(B\mid A)$,
$P(C\mid A)$, plus one cross-family sentence.
\begin{center}\small
\begin{verbatim}
A1: 0.4 <= P(A)        <= 0.6
B1: 0.1 <= P(B | A)    <= 0.3
C1: 0.2 <= P(C | A)    <= 0.5
BC: 0.0 <= P(B and C)  <= 0.05      <- cross-family (scope {B,C} in no family)
\end{verbatim}
\end{center}
The LMC assertion $(B\indep C\mid A)$ is satisfied by \emph{any} product joint, so it is not what
makes this instance loose; the \emph{sentence} \texttt{BC} is. The free product of \eqref{eq:se} can
pick a high-$P(B\mid A)$ vertex and a high-$P(C\mid A)$ vertex at once, reaching $P(B\land C)$ far
above $0.05$ --- a joint in $\DSE$ but not in $\DLCN$.

\subsection{On the \emph{unconditional} marginals CVE is still exact here}
The unconditional bounds happen to be unaffected by the coupling gap on this instance: the exact
$\DLCN$ intervals are
\[
  P(A)=[0.40,0.60],\qquad P(B)=[0.04,0.72],\qquad P(C)=[0.08,0.80],
\]
confirmed by four independent oracles (parametric SLSQP enforcing $B\indep C\mid A$; a
strong-extension vertex oracle; \texttt{CredalJT}/D5 with SCIP; and single-query CVE ---
\texttt{docs/ccte\_exactness.tex}). CVE with \texttt{coupling="off"} returns exactly $P(B)=[0.04,0.72]$
because, for the \emph{unconditional} $P(B)$, the extremal vertex combination does not require
violating \texttt{BC} --- so $\DSE$ and $\DLCN$ agree on this particular marginal, and CVE's exact
$\DSE$ read-off (Proposition~\ref{prop:se}) is the exact LCN bound.

\begin{center}\small
\begin{tabular}{lccc}
\toprule
 & $P(A)$ & $P(B)$ & $P(C)$ \\
\midrule
exact $=$ \CVE{} $=$ oracle $=$ D5 & $[0.40,0.60]$ & $\mathbf{[0.04,0.72]}$ & $[0.08,0.80]$ \\
CCTE \texttt{coupling="off"}          & $[0.40,0.60]$ & $\mathbf{[0.05,0.65]}$ & $[0.08,0.80]$ \\
CCTE \texttt{coupling="cross-family"} (D4) & $[0.40,0.60]$ & $\mathbf{[0.05,0.55]}$ & $[0.08,0.80]$ \\
\bottomrule
\end{tabular}
\end{center}
The contrast with CCTE is the whole point of the per-target schedule. CCTE, using the \emph{same}
vertices, reports $P(B)=[0.05,0.65]$ --- strictly \emph{inside} the true $\DSE$ interval, i.e.\ not
even a valid outer bound --- because its single shared order eliminates $B$ mid-schedule and prunes
(by dominance) the very function that carried the extremal $P(B)$; dominance in a marginalized scope
is not dominance for the eventual read-off. CVE avoids this precisely by recomputing a query-last
order for $B$, so the extremal function survives (Proposition~\ref{prop:se}). This is the
``order-sensitive outer approximation'' the code comment warns about, and it is CVE's per-target
recomputation that immunizes it.

\subsection{Where the gap bites: conditional queries and D4}
The coupling gap $\DSE\supsetneq\DLCN$ does bite on \emph{conditional} queries. The \texttt{.lcn}
header records that $P(B\mid C=1)$ collapses from $[0.0375,0.9625]$ under \texttt{coupling="off"} to
$[0.0375,0.0375]$ under \texttt{coupling="cross-family"}: the free product allows $B$ and $C$ both
near their high vertices given $A$, which the sentence \texttt{BC} forbids. Scheme D4
(Section~\ref{sec:schemes}) removes exactly those vertex \emph{combinations}. In CVE this is two
coordinated changes (\texttt{cve.py:391--455}):
\begin{enumerate}
  \item \textbf{Order augmentation.} The min-fill order is built from the potential scopes
    \emph{plus} the constraint node sets, \texttt{scopes + constraints.constraint\_node\_sets(...)}
    (\texttt{cve.py:408--410}), forcing the coupled atoms $\{A,B,C\}$ to co-occur in a common bucket.
    Without this the constraint scope is never assembled and the filter is inert.
  \item \textbf{Prune-time filtering.} In each bucket, \emph{after} combining and \emph{before}
    marginalizing, \texttt{combined.filter\_infeasible}
    (\texttt{cve.py:441--442}) drops any function whose assembled joint violates a cross-family
    constraint --- checkable only once the bucket scope covers the constraint's atoms, which the
    augmentation guarantees.
\end{enumerate}
D4 is sound for \emph{unconditional} marginals (the extremum sits at a surviving vertex). For
conditional $P(q\mid e)$ the ratio extremum need not sit at a vertex, so D4 is not guaranteed exact;
the reliable exact engine for cross-family coupling is \texttt{CredalJT} (D5). Note also that on the
\emph{unconditional} $P(B)$ above, D4 is \emph{correctly inert}: CVE augments its order and the
filter fires but finds nothing to drop that would change $[0.04,0.72]$ --- consistent with $\DSE$ and
$\DLCN$ already agreeing on that marginal.

\section{The tightening schemes (D1--D5) for CVE}\label{sec:schemes}

The companion note \texttt{docs/tighter\_approximation.tex} identifies two sources of looseness in
the strong-extension bound and five schemes that attack them:
\begin{itemize}
  \item \textbf{Gap A (local-set widening).} Each $K_v$ is computed per family in isolation, so it
    admits distributions no global LCN distribution realizes (it ignores sentences on \emph{other}
    families and cross-family LMC equalities). This widens $\DSE$ before any propagation.
  \item \textbf{Gap B (strong-extension widening).} Even with perfect $K_v$, the free product of
    \eqref{eq:se} lets each family/parent-config pick its vertex independently, violating
    cross-family sentences. This is the $\DSE\supsetneq\DLCN$ gap of Section~\ref{sec:boundary}.
\end{itemize}
Because CVE is exact over $\DSE$ (Proposition~\ref{prop:se}), it reports whatever $\DSE$ it is
given \emph{verbatim} --- so shrinking $\DSE$ toward $\DLCN$ pulls CVE's output inward with no
algorithmic slack of its own to absorb the change. The schemes therefore split into two classes for
CVE exactly as they do for CCTE.

\paragraph{Class 1 --- input-level (D1, D2, D3): free build-stage knobs.} These change only what
\texttt{CredalNetworkVertices.from\_lcn} produces --- tighter local intervals, hence tighter
enumerated vertices --- and CVE consumes them without any code change.
\begin{itemize}
  \item \textbf{D1} (\texttt{method="linear-tight"}): each family's local LP is enriched with every
    in-scope sentence \emph{and} the scope-restricted LMC bilinear equalities, closing Gap~A
    partially. Vacuous $[0,1]$ CPT rows collapse to genuine sub-intervals.
  \item \textbf{D2} (\texttt{merge\_budget=$k$}): merge adjacent families whose combined scope
    $\le k$ before factorization; at $k=1$ no merge, at $k=n$ exact. Larger merged scopes bring
    cross-family constraints in-scope, closing Gap~A and, via the merge, part of Gap~B.
  \item \textbf{D3} (\texttt{solver="scip"}): solve each family's D1 model with SCIP's certified
    global branch-and-bound instead of local ipopt, closing Gap~A \emph{exactly} (and flagging any
    unsound ipopt local optimum). Cost $2^n$ per solve.
\end{itemize}
Being exact over $\DSE$, CVE is the cleanest beneficiary of Class~1: the tighter strong extension is
reported exactly, so D1/D2/D3 pull CVE's interval straight toward $\DLCN$.

\paragraph{Class 2 --- propagation-level coupling (D4, D5): per-algorithm.}
\begin{itemize}
  \item \textbf{D4} (\texttt{coupling="cross-family"} on \texttt{CVE.run}): implemented inside CVE as
    the order augmentation $+$ prune-time filter of Section~\ref{sec:boundary}
    (\texttt{cve.py:391--455}, \texttt{coupling.py:101--250}). It attacks Gap~B by dropping infeasible
    vertex \emph{combinations}. Sound for unconditional marginals; not guaranteed for conditional
    ratios.
  \item \textbf{D5} (\texttt{CredalJT}, \texttt{cn/junction\_nlp.py}): a \emph{separate} engine, not a
    CVE knob. It builds one junction tree, introduces cluster-marginal variables, imposes
    separator-consistency and hosts each LCN constraint on a cluster, and solves one NLP per atom.
    It closes Gap~A $+$ Gap~B \emph{exactly} whenever every sentence is hosted and every LMC is
    hosted-or-separator-entailed (Theorem~1 of \texttt{docs/tighter\_approximation.tex}), at cost
    exponential in \emph{treewidth} rather than $n$. It is sound for both unconditional and
    conditional queries. For cross-family coupling this is the reliable exact engine.
\end{itemize}

\begin{center}\small
\begin{tabular}{lllll}
\toprule
Scheme & Where it lives for CVE & Gap closed & Exact? & Cost \\
\midrule
D1 \texttt{method="linear-tight"} & build (\texttt{from\_lcn}) & A (partial) & no & $\approx 0$ \\
D2 \texttt{merge\_budget=$k$}     & build (\texttt{from\_lcn}) & A $(+\,$B$)$ & at $k=n$ & $2^k$ \\
D3 \texttt{solver="scip"}         & build (\texttt{from\_lcn}) & A (exact)   & Gap~A only & $2^n$/solve \\
D4 \texttt{coupling="cross-family"} & inside \texttt{CVE.run}  & B           & uncond.\ only & med \\
D5 \texttt{CredalJT}              & separate engine           & A $+$ B     & yes (Thm.\ 1) & treewidth \\
\bottomrule
\end{tabular}
\end{center}

\section{Verdict}\label{sec:verdict}

\begin{proposition}[Summary]\label{prop:verdict}
Let $\mathcal L$ be a chain-graph LCN and run \CVE{} with $\epsilon=0$.
\begin{enumerate}
  \item (\emph{Always}) With \texttt{coupling="off"}, CVE returns exactly the strong-extension
    interval $[\lo q_{\mathrm{SE}},\hi q_{\mathrm{SE}}]$, a valid outer bound on $\DLCN$
    (Proposition~\ref{prop:se}).
  \item (\emph{Exact}) If $\mathcal L$ is singly connected and has no cross-family sentence --- Markov
    chain, tree, polytree, or compound-node chain graph such as \texttt{alarm} --- then
    $\DSE=\DLCN$ and CVE is exact for every atom (Proposition~\ref{prop:exact}).
  \item (\emph{Boundary}) A cross-family sentence makes $\DSE\supsetneq\DLCN$ possible. Unconditional
    marginals may still be exact (they are on \texttt{d4\_biting}); conditional queries and the
    general case need D4 (unconditional-sound) or, for a guaranteed exact bound, the D5 engine
    \texttt{CredalJT}.
  \item (\emph{Tightening}) D1/D2/D3 tighten CVE for free at build time and are reported verbatim
    (CVE being exact over $\DSE$); D4 is CVE's built-in coupling filter with order augmentation; D5
    is the separate exact engine.
\end{enumerate}
\end{proposition}

\noindent Compared to its siblings: CVE is exact over $\DSE$ (like CCTE with $\epsilon=0$, unlike
ApproxLP's inner bound and IBP/IJGP's loopy outer slack), and its per-target query-last schedule
makes it \emph{strictly more reliable than CCTE} the moment a cross-family sentence appears --- where
CCTE's single shared order can fall inside the strong-extension interval, CVE always returns it
exactly. The three companion notes cover the neighbors: \texttt{docs/ccte\_exactness.tex} (the
two-pass schedule and its \texttt{d4\_biting} inner error), \texttt{docs/ariel\_exactness.tex}
(structural LMC preservation on singly-connected factor graphs), and
\texttt{docs/tighter\_approximation.tex} (the D1--D5 taxonomy and the D5 exactness theorem).

\end{document}
